\section{Miller sequences and essential incompleteness}\label{sec:miller}

Let $2^{\N}=\{0,1\}^{\N}$ be the set of infinite binary sequences,
with the product topology of the discrete space $\{0,1\}$.
This is Cantor space.
Fix a recursive bijection $\langle-,-\rangle:\N^2\to\N$ and an
acceptable enumeration $(\Phi_e^p)_{e\in\N}$ of partial numerical
oracle computations, where $p\in2^{\N}$. Tuple inputs and outputs use
fixed recursive codings. Let $(q_j)_{j\in\N}$ be a recursive enumeration
of the rationals. An oracle program with index $d\in\N$
\emph{computes a real $r$ from $p$}
if $\Phi_d^p(k)$ is defined for every $k$ and
\[
 |q_{\Phi_d^p(k)}-r|<2^{-k}\qquad(k\in\N).
\]
Give $Q=[0,1]^{\N}$ its product topology and define
\[
 \rho:2^{\N}\longrightarrow Q,\qquad
 \rho(p)(n)=\sum_{j\geq0}p(\langle n,j\rangle)2^{-j-1}.
\]
For each $n$, the subsequence $j\mapsto p(\langle n,j\rangle)$
is a binary expansion of $\rho(p)(n)$.
A \emph{binary representation} of $z\in Q$ is a sequence
$p\in2^{\N}$ with $\rho(p)=z$. Write
$P_z=\rho^{-1}(\{z\})$ for the set of these representations.
The map $\rho$ is continuous and onto, so every $P_z$ is nonempty
and compact.

Miller's construction gives a sequence $a\in Q$ with the following
property: if an oracle program with index $d$ computes the same real $r\in[0,1]$ from
every $p\in P_a$, then
\begin{equation}\label{eq:miller-property}
 a(d)=r.
\end{equation}
Such a sequence is called a \emph{Miller sequence}
\cite[Theorem~2.2]{BauerHanson}.
The index equality is the strengthened conclusion stated after
\cite[Theorem~6.3]{Miller2004}; the formulation with binary
representations is given in
\cite[Section~2.1 and Theorem~2.2]{BauerHanson}.
The rational approximations produced by $d$ may depend on $p$,
while the real they represent is fixed.

\subsection{Computation independent of the binary representation}
\label{sec:miller-category}

\begin{lemma}\label{lem:miller-unanimity}
There is one oracle program which, from $p\in P_z$ and numerical inputs
$e,n$, computes the partial function $\tau_z:\N^2\rightharpoonup\N$
defined by
\[
 \tau_z(e,n)=m
 \quad\Longleftrightarrow\quad
 (\forall q\in P_z)\;\Phi_e^q(n)\mathord\downarrow=m.
\]
Its value and its domain of definition depend only on $z$, not on the
chosen binary representation $p$.
\end{lemma}
\begin{proof}
For a finite binary sequence $\sigma$, let
$[\sigma]=\{p\in2^{\N}:p\mathbin{\upharpoonright}|\sigma|=\sigma\}$
be its cylinder in Cantor space.
The set $\rho([\sigma])$ is a computably specified finite
union of closed dyadic boxes in $Q$, with restrictions on only finitely
many coordinates. For each affected coordinate $n$, complete the
prescribed binary digits to a finite initial segment of
$j\mapsto p(\langle n,j\rangle)$. There are only finitely many such
completions.
From a binary representation $p$ of $z$, strict rational inequalities for the coordinates
of $z$ are semidecidable. We can therefore enumerate all cylinders
$[\sigma]$ for which $z\notin\rho([\sigma])$. Their union is exactly
$2^{\N}\setminus P_z$: if $\rho(q)\ne z$, a sufficiently long initial
segment of $q$ already forces a different value in one coordinate.

For fixed $e,n,m$, let
\[
 H_{e,n,m}=\{q\in2^{\N}:\Phi_e^q(n)\mathord\downarrow=m\}.
\]
Finite halting computations enumerate this open set as a union of
cylinders, uniformly in $e,n,m$. For each $m$, search for a finite cover
of Cantor space by cylinders from $H_{e,n,m}$ and
$2^{\N}\setminus P_z$, dovetailing these searches over $m$. Whether finitely
many cylinders cover Cantor space is decidable. Such a cover exists
exactly when $P_z\subseteq H_{e,n,m}$, by compactness. Since $P_z$ is
nonempty, at most one $m$ satisfies this condition. Thus the search has
the asserted output and definedness, for every binary representation of $z$.
\end{proof}

For the fixed sequence $a$, let $\mathcal E_a$ have the
objects generated from $\N$ by binary Cartesian products and tagged
disjoint unions, with fixed recursive codings. Every such object is
recursively isomorphic to $\N$. A morphism $f:X\rightharpoonup Y$ belongs
to $\mathcal E_a$ when one oracle program computes exactly $f$, including
its domain of definition, for every $p\in P_a$. Equality of morphisms is
equality of partial functions.

\begin{proposition}\label{prop:miller-category}
The category $\mathcal E_a$ is a locally connected recursion category
with ranges, binary restriction coproducts, and finite compatible joins.
It satisfies weak choice, and $\tau_a:\N\times\N\rightharpoonup\N$
is a Turing morphism.
\end{proposition}
\begin{proof}
Composition of the representing programs gives the same composite
partial function for every binary representation of $a$.
The recursive projections, pairings,
tagged injections, and case distinctions give the stated finite
structure. Restriction is computed by running a program for $f(x)$ and,
if it halts, returning $x$. To compute the range partial identity at
$y$, dovetail the computations $f(x)$ over all coded inputs $x$ and
return $y$ when one returns $y$. Although the discovered witness may
depend on the representation, the returned value and the existence of a witness
do not. Thus $\ran(f)$ belongs to $\mathcal E_a$; its leastness and
agreement law follow from those of set-theoretic images.
Dovetailing compatible computations gives their union;
compatibility ensures that the output is independent of which
computation finishes first. These operations also give zero maps and
the prodominical law $f\times0=0$.

For every input $x$, the singleton partial identity $1_{\{x\}}$ is
recursive. It is connected, and these partial identities detect
inequality of partial functions, so the category is locally connected.
A partial monomorphism here is injective on its domain of definition: a failure of
injectivity would be detected by two recursive constant maps. For an
injective morphism $f$, the same preimage search returns the unique $x$
with $f(x)=y$. Its output is consequently independent of the representation and
gives the partial inverse. This proves weak choice.

Lemma~\ref{lem:miller-unanimity} puts $\tau_a$ in $\mathcal E_a$.
Let $f:X\times\N\rightharpoonup\N$ be represented by a program with
index $d$. The parameter theorem for oracle programs gives an ordinary
total recursive map $c:X\to\N$ with
\[
 \Phi_{c(x)}^p(n)\simeq\Phi_d^p(x,n).
\]
Here $\simeq$ includes simultaneous undefinedness.
If $f(x,n)=m$, every $p\in P_a$ gives this same output; if $f(x,n)$ is
undefined, every $p\in P_a$ gives an undefined computation. The definition
of $\tau_a$ therefore gives
$f=\tau_a(c\times1_{\N})$. The map $c$ is a total source morphism.
Finally, the chosen recursive codings make the category isotypical.
\end{proof}

For $p\in P_a$, let $\mathrm{Rec}^{p}$ be the recursion category on
the same numerical objects, with partial $p$-recursive functions as
morphisms. The identity interpretation of partial functions defines a
faithful range functor
\[
 J_p:\mathcal E_a\longrightarrow\mathrm{Rec}^{p}.
\]
It also preserves restriction products and finite compatible joins.
Write $\mathcal K_p=\operatorname{Split}(\mathrm{Rec}^{p})$ and let
$I_p:\mathrm{Rec}^{p}\to\mathcal K_p$ be the full inclusion.
This inclusion is a range functor.
A concrete presentation of $\mathcal K_p$ has $p$-c.e. subsets of the
numerical objects as objects and partial $p$-recursive functions between
these carriers as morphisms. A restriction idempotent splits through its
$p$-c.e. domain of definition.
Singleton restrictions prove local connectedness, and preimage search
for an injective map proves weak choice. Ranges are the partial identities
on images, and inverse images preserve unions of domains of definition.
Thus $\mathcal K_p$ satisfies the target
hypotheses of Theorem~\ref{thm:local-lifting}.

If $c,c'$ are complementary restriction idempotents on a numerical
object $X$, run the membership computations for their domains of
definition in parallel and return the opposite tag when one succeeds.
This gives the total source map
\[
 \ell_c=(i_1c)\vee(i_0c'):X\longrightarrow X+X.
\]
Its image under $I_pJ_p$ is the required section, and
$(I_pJ_p)(1_{\N})\ne0$. Thus each composite
$I_pJ_p$ satisfies the hypotheses of Theorem~\ref{thm:local-lifting}.
Using the recursive bijection $b:\N\to\N+\N$ which assigns the tag
of $m$ according to its parity, that theorem gives inseparable
restriction idempotents that are the partial identities on the sets
\begin{equation}\label{eq:miller-diagonal-pair}
 A_i=\{e:\tau_a(e,e)\mathord\downarrow=m
                   \text{ for some }m\equiv i\pmod2\},
 \qquad i=0,1.
\end{equation}

\subsection{The index sets below and above \texorpdfstring{$1/2$}{1/2}}
\label{sec:miller-cuts}

Define
\[
 L_a=\{n:a(n)<\tfrac12\},\qquad
 R_a=\{n:a(n)>\tfrac12\}.
\]
Their partial identities belong to $\mathcal E_a$, since strict
inequalities can be recognized from approximations obtained from any
binary representation of $a$.

\begin{lemma}\label{lem:miller-cut-reduction}
There is an ordinary total recursive map $h:\N\to\N$ such that
\[
 h^{-1}(L_a)=A_0,\qquad h^{-1}(R_a)=A_1.
\]
\end{lemma}
\begin{proof}
For $e,m,t\in\N$, let $H_{e,m,t}$ be the finite union of halting
cylinders enumerated by stage $t$ for the computation
$\Phi_e^q(e)\mathord\downarrow=m$. Choose these finite unions increasing
in $t$, and put
\[
 K_{e,m,t}=\rho(2^{\N}\setminus H_{e,m,t})\subseteq Q.
\]
Each $K_{e,m,t}$ is a computably specified finite union of closed dyadic
boxes. Compactness of $P_z$ gives
\begin{align*}
 U_e^i
 &:=\{z\in Q:\tau_z(e,e)\mathord\downarrow=m
                      \text{ for some }m\equiv i\pmod2\}\\
 &=\bigcup_{\substack{m\equiv i\pmod2\\t\in\N}}
       (Q\setminus K_{e,m,t}).
\end{align*}
In particular, $U_e^0$ and $U_e^1$ are disjoint effectively open sets.

The product topology on $Q$ is induced by the metric
\[
 d:Q\times Q\to[0,1],\qquad
 d(z,w)=\sum_{n\geq0}2^{-n-1}|z(n)-w(n)|.
\]
For a nonempty box restricting finitely many coordinates to closed
rational intervals $I_n$, the distance from $z$ to the box is the
finite sum
$\sum_n2^{-n-1}\operatorname{dist}(z(n),I_n)$ over those coordinates.
Distance to a nonempty finite union is the minimum of these distances.
Define $\delta_{e,m,t}:Q\to[0,1]$ by
$\delta_{e,m,t}(z)=d(z,K_{e,m,t})$ when $K_{e,m,t}$ is nonempty, and
$\delta_{e,m,t}(z)=1$ when it is empty. These functions are uniformly
computable and continuous, and
\[
 \delta_{e,m,t}(z)>0\quad\Longleftrightarrow\quad
 z\notin K_{e,m,t}.
\]
For each $i$, take a fixed recursive enumeration
$(m_j,t_j)_{j\geq0}$ of all pairs with $m_j\equiv i\pmod2$, and define
\[
 g_e^i:Q\to[0,1],\qquad
 g_e^i(z)=\sum_{j\geq0}2^{-j-1}\delta_{e,m_j,t_j}(z).
\]
The tail bound is uniform, so $g_e^i$ is computable and continuous,
with $0\leq g_e^i\leq1$ and $g_e^i(z)>0$ exactly on $U_e^i$.
Consequently the computable continuous function
\[
 r_e:Q\to[1/4,3/4],\qquad
 r_e(z)=\frac12+\frac{g_e^1(z)-g_e^0(z)}4
\]
is below $1/2$ exactly on $U_e^0$, and
is above $1/2$ exactly on $U_e^1$.

The preceding construction gives one oracle program computing total
numerical functions $\Psi^p:\N^2\to\N$, for $p\in2^{\N}$.
Its output $\Psi^p(e,k)$ is an index in the rational enumeration,
with
\[
 \bigl|q_{\Psi^p(e,k)}-r_e(\rho(p))\bigr|<2^{-k}
 \qquad(e,k\in\N,\ p\in2^{\N}).
\]
The oracle parameter theorem gives a total recursive $h$ such that
$\Phi_{h(e)}^p(k)=\Psi^p(e,k)$. For $p\in P_a$ these computations
represent the same real $r_e(a)$, so the Miller property gives
\[
 a(h(e))=r_e(a).
\]
Together with $a\in U_e^i\Longleftrightarrow e\in A_i$, this proves
the asserted inverse-image equalities.
\end{proof}

\begin{corollary}\label{cor:miller-cuts-inseparable}
The restriction idempotents $1_{L_a},1_{R_a}:\N\rightharpoonup\N$ are inseparable in
$\mathcal E_a$. Nevertheless, every $p\in P_a$ computes a separator
of the underlying sets $L_a,R_a$.
\end{corollary}
\begin{proof}
If complementary restriction idempotents in $\mathcal E_a$ separated
$1_{L_a},1_{R_a}$, their inverse images under $h$ would separate
$1_{A_0},1_{A_1}$, contrary to
Theorem~\ref{thm:local-lifting} as applied above.
For a fixed binary representation $p\in P_a$, the set
\[
 C_p=\{n:p(\langle n,0\rangle)=0\}
\]
is $p$-decidable and satisfies $L_a\subseteq C_p$ and
$C_p\cap R_a=\varnothing$. At $a(n)=1/2$, either first binary digit
is allowed. Thus the single program reading that digit produces a
separator from every binary representation of $a$, but the set it decides
cannot be independent of the representation.
\end{proof}

If $J_p\tau_a$ were a Turing morphism, Theorem~\ref{thm:local-lifting},
applied to $I_p:\mathrm{Rec}^{p}\to\mathcal K_p$ and its
$p$-recursive sections associated with complementary restriction
idempotents, would make $1_{A_0},1_{A_1}$ inseparable in
$\mathrm{Rec}^{p}$. The $p$-decidable set $h^{-1}(C_p)$ separates
$A_0,A_1$. Hence $J_p\tau_a$ is not a Turing
morphism, although $\mathrm{Rec}^{p}$ has a Turing morphism.


\subsection{A theory of real closed fields}
\label{def:miller-splitting}

Call a subset $D$ of a numerical object \emph{$\mathcal E_a$-semidecidable}
when its partial identity $1_D$ is a morphism of $\mathcal E_a$.
We use the concrete presentation of $\operatorname{Split}(\mathcal E_a)$
in which such a subset $D\subseteq X$ represents $(X,1_D)$.
For another $\mathcal E_a$-semidecidable subset $E\subseteq Y$,
a morphism $D\rightharpoonup E$ is represented by
a morphism $f:X\rightharpoonup Y$ of $\mathcal E_a$ satisfying
$f=1_Ef1_D$; it is total when $\bar f=1_D$.
All code operations below are restricted to their indicated carriers.

Put $\mathcal D_a=\operatorname{Tot}(\operatorname{Split}(\mathcal E_a))$.
The following realization permits us to apply the categorical theory
construction of Appendix~\ref{app:syntax-theories}.

\begin{proposition}\label{prop:miller-theory-realization}
The category $\mathcal D_a$ is regular, with finite coproducts that
are disjoint and stable under pullback. Every monomorphism is an
admissible domain monomorphism, and there is a restriction equivalence
\[
 \operatorname{Par}(\mathcal D_a,\operatorname{Mono})
       \simeq\operatorname{Split}(\mathcal E_a).
\]
For $V=\N$ and a recursive signature of bounded finite arity, the
syntax and list algebras of Definition~\ref{def:categorical-syntax}
exist in $\mathcal D_a$ and are stable under parameters.
For every axiom subobject of the sentence object, the inference
endofunctor of Theorem~\ref{thm:first-order-theory} has an initial
algebra of finite derivations.
\end{proposition}
\begin{proof}
A singleton is terminal. Finite products and equalizers are
$\mathcal E_a$-semidecidable carriers, using parallel membership
tests and numerical equality. The empty carrier and tagged unions
give finite coproducts.
For a morphism $f:C\to D+E$, running $f$ and testing its output tag
gives two semidecidable inverse-image carriers. Their coproduct
maps isomorphically to $C$, with inverse given by the same tag test.
The injections have disjoint images, so these coproducts are disjoint
and stable under pullback.

For $f:C\to D$, dovetail its computations over all numerical inputs
to semidecide its image. A surjection $e:C\to D$ is the coequalizer
of its kernel pair. Indeed, for $h:C\to Z$ constant on the fibers
of $e$, search on input $y$ for $x$ with $e(x)=y$ and return $h(x)$.
This computation has domain exactly $D$. The witness found may
depend on the binary representation of $a$, but its $h$-value does
not. It therefore defines the unique descended morphism $D\to Z$.
Every morphism factors through its image in this way, and pullbacks
preserve surjections, proving regularity. Conversely, a regular
epimorphism is surjective, since its image factorization and
extremality force the image inclusion to be an isomorphism.

Singleton maps detect noninjectivity, so the monomorphisms are
exactly the injections. For an injection $m:C\to D$, preimage
search returns the unique $x$ with $m(x)=y$. Its domain is $m(C)$,
and both its domain and its value are independent of the binary
representation. This gives a partial inverse $m^\dagger$ in the
restriction splitting. Hence every monomorphism is an admissible
domain inclusion up to isomorphism. A span
$C\xleftarrow{m}B\xrightarrow{g}D$ represents $gm^\dagger$;
conversely, every partial morphism is total on its semidecidable
domain. These assignments respect composition by pullback and give
the displayed restriction equivalence.

Numerical equality makes $V=\N$ decidable. Encode terms and formulas
as finite labeled ordered trees with distinct constructor tags,
and lists of variable names as finite sequences. Recursive parsing
and the recursive symbol carriers make these carriers and their
constructors recursive. Fix an arbitrary object $P$ of $\mathcal D_a$
and a $K_P$-algebra on two objects over $P$. On input $(w,t)$ or
$(w,\varphi)$, first semidecide $w\in P$ and check the tree, then
evaluate it from its leaves to its root using the algebra operations
in the fiber over $w$. Each operation preserves that fiber.
Induction on the tree proves termination and gives the same result
for every $p\in P_a$, because each algebra operation has a common
domain and value. Thus evaluation defines morphisms from
$P\times\mathsf{Tm}$ and $P\times\mathsf{Fm}$ with exactly these
source domains. They satisfy the algebra equations, and induction
proves uniqueness. This proves syntax initiality in every slice
over $P$. Evaluation of finite lists, retaining the same parameter
$w$, proves the corresponding assertion for the list algebra.

The free-variable and substitution folds of the appendix consequently
exist in $\mathcal D_a$. In this realization, the free-variable fold
computes finite lists of occurrences. The single-variable
substitution fold computes either failure or a resulting formula:
it compares the bound and substituted variables and uses the finite
membership tests in~\eqref{eq:quantifier-fold}. These computations
are ordinary recursive functions. The pullback in
Proposition~\ref{prop:categorical-substitution} therefore gives the
recursive free-for domain and substitution on that domain. The
sentence carrier, where the free-variable list is empty, is recursive.

Represent an axiom subobject by its semidecidable image in the
sentence carrier. Use the seven logical schemes, modus ponens,
and generalization of Appendix~\ref{app:first-order-derivations}.
A finite derivation tree records its rule instance and ordered
children at every node. Tree shape, logical instance conditions,
and agreement of child conclusions with the required premises are
recursive tests. Run these tests together with the finitely many
nonlogical axiom membership tests. They semidecide a carrier
$P_A$ with the same membership outcome for every $p\in P_a$.
Reading the root gives a total conclusion morphism
$c_A:P_A\to\mathsf{Fm}$, and adjoining a labeled root gives an
$H_A$-algebra.

For any $H_A$-algebra, first semidecide that the input is a valid
derivation, then evaluate it from its leaves to its root. Induction
shows that each intermediate value has the prescribed conclusion,
so the next algebra operation is applied on its domain. The
evaluation terminates on every valid derivation. The operations'
common domains and values make its domain and result independent
of $p\in P_a$, so it is a morphism of $\mathcal D_a$ over
$\mathsf{Fm}$. Its recursive equations make it an algebra morphism;
induction on derivations proves uniqueness. Thus the derivation
algebra is initial.
\end{proof}

Consider the language of ordered fields with additional constants
$(c_n)_{n\in\N}$. Its symbol carriers are recursive and its arities
are at most two, so Proposition~\ref{prop:miller-theory-realization}
applies. Let $T_a$ consist of a recursive axiomatization of
real closed ordered fields together with the axioms
\begin{equation}\label{eq:miller-theory}
 q<c_n\quad\text{if }q<a(n),\qquad
 c_n<q\quad\text{if }a(n)<q,
 \qquad(q\in\mathbb Q, n\in\N).
\end{equation}
Rational bounds abbreviate inequalities with integer coefficients.
The real field, with $c_n$ interpreted as $a(n)$, is a model of $T_a$.
Its axiom set is $\mathcal E_a$-semidecidable: the field axioms are
recursive, and each strict rational inequality can be recognized from
every binary representation of $a$ with the same membership outcome.

\begin{theorem}\label{thm:miller-incompleteness}
Let $a$ be a Miller sequence with property~\eqref{eq:miller-property},
and let $T_a$ be the theory~\eqref{eq:miller-theory}.
For every consistent extension of $T_a$ with an
$\mathcal E_a$-semidecidable axiom set, the partial identities on its
theorem and refutation sets are inseparable in
$\operatorname{Split}(\mathcal E_a)$.
In particular, every such extension is incomplete.
\end{theorem}
\begin{proof}
The set $S$ of sentence codes is recursive. It therefore defines an
object of $\operatorname{Split}(\mathcal E_a)$. Negation and insertion
of a constant into an inequality give ordinary recursive total morphisms
\[
 \nu:S\to S,\qquad
 \nu(\ulcorner\varphi\urcorner)=\ulcorner\neg\varphi\urcorner,
 \qquad
 \theta:\N\to S,\qquad
 \theta(n)=\ulcorner c_n<1/2\urcorner.
\]
For a theory $T$ with an $\mathcal E_a$-semidecidable axiom set,
let $A_T\hookrightarrow S$ be its axiom subobject. By
Proposition~\ref{prop:miller-theory-realization}, the inference
endofunctor $H_{A_T}$ has an initial derivation algebra. Pull its
conclusion morphism back to $S$, obtaining
\[
 c_T:\mathrm{Prf}_T\to S.
\]
Theorem~\ref{thm:first-order-theory} identifies its image with the
sentence theory generated by $A_T$. The realization proposition
makes this image monomorphism an admissible domain.
Lemma~\ref{lem:range-regular-image} therefore identifies
$t_T=\ran(c_T)$ with the partial identity on the theorem set in
$\operatorname{Split}(\mathcal E_a)$.
The inverse image $\nu^*t_T$ is the partial identity on the set of
sentence codes whose negations are provable in $T$.

For $n\in L_a$, the sentence $c_n<1/2$ is an axiom of $T_a$.
For $n\in R_a$, the axiom $1/2<c_n$ and the ordered-field axioms
prove its negation. Consequently
\[
 1_{L_a}\le\theta^*t_{T_a},\qquad
 1_{R_a}\le\theta^*\nu^*t_{T_a}.
\]
By Corollary~\ref{cor:miller-cuts-inseparable}, the restriction idempotents
$1_{L_a},1_{R_a}$ are inseparable; the full inclusion into the restriction
splitting preserves this property. Joins of restriction idempotents in
the splitting are partial identities on unions, and their inverse images
are partial identities on the corresponding set preimages.
For a consistent extension $T\supseteq T_a$, leastness in
Theorem~\ref{thm:first-order-theory} gives $t_{T_a}\le t_T$,
while consistency gives $t_T\nu^*t_T=0_S$.
Corollary~\ref{cor:domain-incompleteness} therefore makes
$t_T,\nu^*t_T$ inseparable and gives
$t_T\vee\nu^*t_T\ne1_S$. This is incompleteness of $T$.
\end{proof}
